% Options for packages loaded elsewhere
\PassOptionsToPackage{unicode}{hyperref}
\PassOptionsToPackage{hyphens}{url}
\documentclass[
]{article}
\usepackage{xcolor}
\usepackage{amsmath,amssymb}
\setcounter{secnumdepth}{-\maxdimen} % remove section numbering
\usepackage{iftex}
\ifPDFTeX
  \usepackage[T1]{fontenc}
  \usepackage[utf8]{inputenc}
  \usepackage{textcomp} % provide euro and other symbols
\else % if luatex or xetex
  \usepackage{unicode-math} % this also loads fontspec
  \defaultfontfeatures{Scale=MatchLowercase}
  \defaultfontfeatures[\rmfamily]{Ligatures=TeX,Scale=1}
\fi
\usepackage{lmodern}
\ifPDFTeX\else
  % xetex/luatex font selection
\fi
% Use upquote if available, for straight quotes in verbatim environments
\IfFileExists{upquote.sty}{\usepackage{upquote}}{}
\IfFileExists{microtype.sty}{% use microtype if available
  \usepackage[]{microtype}
  \UseMicrotypeSet[protrusion]{basicmath} % disable protrusion for tt fonts
}{}
\makeatletter
\@ifundefined{KOMAClassName}{% if non-KOMA class
  \IfFileExists{parskip.sty}{%
    \usepackage{parskip}
  }{% else
    \setlength{\parindent}{0pt}
    \setlength{\parskip}{6pt plus 2pt minus 1pt}}
}{% if KOMA class
  \KOMAoptions{parskip=half}}
\makeatother
\usepackage{longtable,booktabs,array}
\newcounter{none} % for unnumbered tables
\usepackage{calc} % for calculating minipage widths
% Correct order of tables after \paragraph or \subparagraph
\usepackage{etoolbox}
\makeatletter
\patchcmd\longtable{\par}{\if@noskipsec\mbox{}\fi\par}{}{}
\makeatother
% Allow footnotes in longtable head/foot
\IfFileExists{footnotehyper.sty}{\usepackage{footnotehyper}}{\usepackage{footnote}}
\makesavenoteenv{longtable}
\setlength{\emergencystretch}{3em} % prevent overfull lines
\providecommand{\tightlist}{%
  \setlength{\itemsep}{0pt}\setlength{\parskip}{0pt}}
\usepackage{bookmark}
\IfFileExists{xurl.sty}{\usepackage{xurl}}{} % add URL line breaks if available
\urlstyle{same}
\hypersetup{
  hidelinks,
  pdfcreator={LaTeX via pandoc}}

\author{}
\date{}

\begin{document}

\section{References}\label{references}

Working bibliography for THEORY.md and THEORY\_FEP.md. Citation keys are
in \texttt{{[}AuthorYear{]}} format for inline placement.

\begin{center}\rule{0.5\linewidth}{0.5pt}\end{center}

\subsection{Tier 1 --- Must cite before
submission}\label{tier-1-must-cite-before-submission}

\subsubsection{{[}Friston2010{]}}\label{friston2010}

Friston, K. (2010). The free-energy principle: a unified brain theory?
\emph{Nature Reviews Neuroscience}, 11(2), 127--138.
https://doi.org/10.1038/nrn2787

\textbf{Supports:} Abstract, Section 4 (Theorem 2a note), THEORY\_FEP.md
throughout. \textbf{Stance:} Extension and partial correction. FEP
derives that intelligent systems minimize free energy (equivalent to
U-reduction) but does not derive the minimization imperative from
selection pressure, does not constrain the generative model via K/F
inseparability, assumes top-down precision allocation (which Theorem 2a
shows is self-undermining under C\_n constraints), and has no account of
collective intelligence.

\subsubsection{{[}Friston2021{]}}\label{friston2021}

Friston, K., Da Costa, L., Hafner, D., Hesp, C., \& Parr, T. (2021).
Sophisticated inference. \emph{Neural Computation}, 33(3), 713--763.
https://doi.org/10.1162/neco\_a\_01351

\textbf{Supports:} Section 11 Related Work (FEP extension 4 ---
collective intelligence); THEORY\_FEP.md Section 6 (Extension 4,
multi-agent FEP). \textbf{Stance:} Engaged directly. Sophisticated
inference extends active inference to recursive, multi-step planning and
multi-agent coordination. Extension 4 shows what this framework lacks: a
formal fidelity threshold (\(\lambda_{min}\)), provenance as the
mechanism maintaining it, and truth as the limit of collective
action-validation.

\subsubsection{{[}Friston2017{]}}\label{friston2017}

Friston, K., FitzGerald, T., Rigoli, F., Schwartenbeck, P., \& Pezzulo,
G. (2017). Active inference: A process theory. \emph{Neural
Computation}, 29(1), 1--49. https://doi.org/10.1162/NECO\_a\_00912

\textbf{Supports:} THEORY\_FEP.md Section on precision hierarchy;
Theorem 2a note on top-down allocation. \textbf{Stance:} Correction.
Active inference formalizes the precision-weighted prediction error as
the attention allocation mechanism. This is the specific target of the
allocation tax argument and the escalation critique.

\subsubsection{{[}Shannon1948{]}}\label{shannon1948}

Shannon, C. E. (1948). A mathematical theory of communication.
\emph{Bell System Technical Journal}, 27(3), 379--423.
https://doi.org/10.1002/j.1538-7305.1948.tb01338.x

\textbf{Supports:} Definition 5 (U = H(w\textbar K)), Definition 3b (G =
I(f;w\textbar K)), all information-theoretic grounding throughout.
\textbf{Stance:} Foundation. H(w\textbar K) is Shannon conditional
entropy; I(f;w\textbar K) is Shannon mutual information. The paper's
uncertainty measure is this quantity exactly.

\subsubsection{{[}Tarski1936{]}}\label{tarski1936}

Tarski, A. (1936). The concept of truth in formalized languages. In J.
H. Woodger (Trans.), \emph{Logic, Semantics, Metamathematics}
(pp.~152--278). Oxford University Press (1956).

\textbf{Supports:} Theorem 10 (Epistemic Ceiling); Definition 16a (Truth
as fixed point). \textbf{Stance:} Formal grounding. Tarski's
undefinability theorem establishes that truth in a formal system cannot
be defined within that system. Theorem 10 is the information-theoretic
operational form of this result: the confidence regress is non-closeable
from within, which is the epistemic consequence of Tarski's theorem
applied to the framework's truth definition.

\subsubsection{{[}Peirce1877{]}}\label{peirce1877}

Peirce, C. S. (1877). The fixation of belief. \emph{Popular Science
Monthly}, 12, 1--15.

\subsubsection{{[}Peirce1878{]}}\label{peirce1878}

Peirce, C. S. (1878). How to make our ideas clear. \emph{Popular Science
Monthly}, 12, 286--302.

\textbf{Supports:} Definition 16 (Truth as limit of collective
action-validation). \textbf{Stance:} Correspondence. Peirce defines
truth as ``the opinion which is fated to be ultimately agreed to by all
who investigate.'' Definition 16 formalizes this as an
information-theoretic limit: the fraction of independent entities whose
action-validation of p converges as n → ∞. The grounding in
action-validation (not mere agreement) and the explicit convergence
condition are the formal additions this paper makes.

\begin{center}\rule{0.5\linewidth}{0.5pt}\end{center}

\subsection{Tier 2 --- Should cite}\label{tier-2-should-cite}

\subsubsection{{[}Schmidhuber2010{]}}\label{schmidhuber2010}

Schmidhuber, J. (2010). Formal theory of creativity, fun, and intrinsic
motivation (1990--2010). \emph{IEEE Transactions on Autonomous Mental
Development}, 2(3), 230--247. https://doi.org/10.1109/TAMD.2010.2056368

\textbf{Supports:} Definition 6 (Agency A as intrinsic drive),
Definition 11a (η\_M as learning efficiency), Section 1 (intelligence as
driven capacity to improve). \textbf{Stance:} Correspondence and
extension. Schmidhuber formalizes curiosity as compression progress ---
the rate at which a learning algorithm improves its world model. This is
structurally η\_M. The difference: (1) Schmidhuber's drive is about
compression rate, not uncertainty reduction; (2) this paper separates
drive (A) from mechanism (η\_M) as formally independent factors; (3) A
is grounded in survival pressure (Theorem 3) rather than defined as a
design objective.

\subsubsection{{[}Gibson1979{]}}\label{gibson1979}

Gibson, J. J. (1979). \emph{The Ecological Approach to Visual
Perception}. Houghton Mifflin.

\textbf{Supports:} Section 3 (K/F Inseparability), Corollary 1 (Indexed
Knowledge), Definition 3 (Informed Action Space). \textbf{Stance:}
Correspondence and formalization. Gibson's affordances are the action
possibilities a world offers relative to an organism's capacities --- K
indexed by F under a different name. Theorem 1 provides what Gibson's
account lacks: a formal selection criterion (the asymmetric retention
criterion from survival pressure) for which affordances are retained and
which are pruned.

\subsubsection{{[}Baars1988{]}}\label{baars1988}

Baars, B. J. (1988). \emph{A Cognitive Theory of Consciousness}.
Cambridge University Press.

\textbf{Supports:} Definition 8 (Encoding Hierarchy), Theorem 2
(Escalation), Corollary 3 (Context Window as Frontier), Theorem 2a
(Top-Down Allocation Degrades Inference Quality). \textbf{Stance:}
Correspondence and formalization. Global Workspace Theory (GWT) holds
that local processors handle routine computation; the global workspace
broadcasts only when local processors fail. This is the escalation
architecture in cognitive science language. The paper's contribution:
(1) derives the escalation criterion from U-reduction and survival
pressure rather than assuming it architecturally; (2) makes the
hierarchy continuous (Definition 8) rather than binary (local
vs.~global); (3) shows via Theorem 2a why top-down allocation from the
global workspace is inefficient.

\subsubsection{{[}Dehaene2011{]}}\label{dehaene2011}

Dehaene, S., \& Changeux, J.-P. (2011). Experimental and theoretical
approaches to conscious processing. \emph{Neuron}, 70(2), 200--227.
https://doi.org/10.1016/j.neuron.2011.03.018

\textbf{Supports:} Same as {[}Baars1988{]}; use alongside or as
alternative. \textbf{Stance:} Same as {[}Baars1988{]}. Dehaene \&
Changeux provide the neural implementation evidence for GWT, making this
the more empirically grounded citation for the escalation architecture
claim.

\subsubsection{{[}Kahneman2011{]}}\label{kahneman2011}

Kahneman, D. (2011). \emph{Thinking, Fast and Slow}. Farrar, Straus and
Giroux.

\textbf{Supports:} Definition 8 (Encoding Hierarchy), Section 1
(continuity claim). \textbf{Stance:} Discrete precursor that the
continuity principle corrects. System 1 (automatic) and System 2
(deliberate) are the popular formulation of the encoding hierarchy ---
but as a binary. The paper's contribution is making this continuous
(Definition 8, L(c,r) parameterization) and deriving the boundary
criterion from U-reduction cost rather than asserting it.

\subsubsection{{[}Chalmers1995{]}}\label{chalmers1995}

Chalmers, D. J. (1995). Facing up to the problem of consciousness.
\emph{Journal of Consciousness Studies}, 2(3), 200--219.

\textbf{Supports:} Definition 14 (Sentience), note on usage.
\textbf{Stance:} Explicit disclaimer. The paper uses ``sentience'' as a
functional term (degree to which A drives knowing beyond survival
justification) and explicitly makes no claim about phenomenal
consciousness. Chalmers' hard problem is the reason for that disclaimer.
Cite in the Definition 14 note: the functional condition defined here is
not offered as a solution to the hard problem; whether it is sufficient,
necessary, or orthogonal to phenomenal experience is left open.

\subsubsection{{[}Tishby2011{]}}\label{tishby2011}

Tishby, N., \& Polani, D. (2011). Information theory of decisions and
actions. In A. Cutsuridis, A. Hussain, \& J. G. Taylor (Eds.),
\emph{Perception-Action Cycle: Models, Architectures, and Hardware}
(pp.~601--636). Springer. https://doi.org/10.1007/978-1-4419-1452-1\_19

\textbf{Supports:} Section 3 (K/F Inseparability), Definition 3
(Informed Action Space). \textbf{Stance:} Correspondence. The
information bottleneck principle applied to action selection: an agent
retains only information useful for action. Theorem 1's asymmetric
retention criterion is a formal extension grounded in survival pressure,
which Tishby \& Polani's framework treats as a design parameter rather
than a derived result.

\subsubsection{{[}Thrun1998{]}}\label{thrun1998}

Thrun, S., \& Pratt, L. (1998). Learning to learn: Introduction and
overview. In S. Thrun \& L. Pratt (Eds.), \emph{Learning to Learn}
(pp.~3--17). Springer. https://doi.org/10.1007/978-1-4615-5529-2\_1

\textbf{Supports:} Definition 11 (Higher-Order Function Space M),
Definition 11a (η\_M), Definition 7 (I(E) = A·η\_M). \textbf{Stance:}
Correspondence. Meta-learning is the CS field that formalizes the rate
at which a system improves its own learning given experience --- η\_M.
The paper connects meta-learning rate to survival pressure (Theorem 4
rigor thresholds) rather than treating it as a design parameter, and
separates drive (A) from mechanism (η\_M) as formally distinct
quantities.

\subsubsection{{[}Schrodinger1944{]}}\label{schrodinger1944}

Schrödinger, E. (1944). \emph{What is Life? The Physical Aspect of the
Living Cell}. Cambridge University Press.

\textbf{Supports:} Section 5 thermodynamic commentary, Open Question 1
(origin of A), Open Question 9 (DNA as encoding of M). \textbf{Stance:}
Foundation. Schrödinger's ``negative entropy'' argument --- that living
systems maintain local order against thermodynamic pressure --- is the
physical grounding for the claim that U\_lethal connects to
thermodynamic thresholds and that A is thermodynamically instantiated.

\begin{center}\rule{0.5\linewidth}{0.5pt}\end{center}

\subsection{Tier 3 --- Cite for
completeness}\label{tier-3-cite-for-completeness}

\subsubsection{{[}Pearl1988{]}}\label{pearl1988}

Pearl, J. (1988). \emph{Probabilistic Reasoning in Intelligent Systems:
Networks of Plausible Inference}. Morgan Kaufmann.

\textbf{Supports:} Definition 5 (H(w\textbar K) --- how K as a directed
graph induces P(w\textbar K)). \textbf{Stance:} Foundation. Bayesian
networks are the standard formalism for conditioning on graph-structured
knowledge; Definition 5's interpretation of K as a Bayesian network
where nodes are propositions and edges are conditional dependencies uses
Pearl's framework directly.

\subsubsection{{[}Cover2006{]}}\label{cover2006}

Cover, T. M., \& Thomas, J. A. (2006). \emph{Elements of Information
Theory} (2nd ed.). Wiley-Interscience.

\textbf{Supports:} Definition 5 (H(w\textbar K)), Definition 3b
(I(f;w\textbar K)), Lemma 1 proof (monotonicity of conditional entropy).
\textbf{Stance:} Standard reference. Cite for the monotonicity of
conditional entropy (Lemma 1 proof), the data processing inequality
(Definition 3b G≥0 claim), and as the textbook reference for all
information-theoretic identities used.

\subsubsection{{[}Prigogine1984{]}}\label{prigogine1984}

Prigogine, I., \& Stengers, I. (1984). \emph{Order Out of Chaos: Man's
New Dialogue with Nature}. Bantam Books.

\textbf{Supports:} Open Question 1 (thermodynamic instantiation of A),
Open Question 9 (DNA and dissipative structures). \textbf{Stance:}
Adjacent framework. Dissipative structures --- self-organizing systems
that maintain local order by dissipating entropy --- are the physical
precursor of A in primitive form. Already cited in text; add full
reference.

\subsubsection{{[}England2013{]}}\label{england2013}

England, J. L. (2013). Statistical physics of self-replication.
\emph{Journal of Chemical Physics}, 139(12), 121923.
https://doi.org/10.1063/1.4818538

\textbf{Supports:} Open Question 1 (thermodynamic instantiation of A),
Open Question 9 (DNA as encoding of M, convergence from thermodynamics).
\textbf{Stance:} Adjacent framework. England's dissipative adaptation
shows that self-replicating structures are thermodynamically favored in
driven systems. Already cited in text; add full reference.

\subsubsection{{[}MaynardSmith1995{]}}\label{maynardsmith1995}

Maynard Smith, J., \& Szathmáry, E. (1995). \emph{The Major Transitions
in Evolution}. W. H. Freeman.

\textbf{Supports:} Open Question 9 (DNA as encoding of M, evolutionary
encoding hierarchy), Theorem 3a (multi-timescale gradient).
\textbf{Stance:} Correspondence. Major transitions in evolution (e.g.,
the origin of the genetic code, eukaryotes, multicellularity) correspond
to shifts in what level of the encoding hierarchy holds the primary
gradient signal. Each transition is a change in where M operates and at
what rigor threshold.

\subsubsection{{[}Hinton1987{]}}\label{hinton1987}

Hinton, G. E., \& Nowlan, S. J. (1987). How learning can guide
evolution. \emph{Complex Systems}, 1(3), 495--502.

\textbf{Supports:} Theorem 3a (multi-timescale gradient, η\_evo vs
η\_ind), Open Question 9. \textbf{Stance:} Correspondence. The Baldwin
effect --- learned adaptations can guide genetic assimilation --- is the
empirical phenomenon corresponding to the multi-timescale gradient.
Knowledge that first appears at L\_n (individual learning) can
eventually meet θ\_0 and crystallize at L\_0 (genetic encoding) over
evolutionary timescales.

\subsubsection{{[}Sutton2018{]}}\label{sutton2018}

Sutton, R. S., \& Barto, A. G. (2018). \emph{Reinforcement Learning: An
Introduction} (2nd ed.). MIT Press.

\textbf{Supports:} Theorem 3 (gradient descent on U), Theorem 3a
(multi-timescale gradient). \textbf{Stance:} Correspondence.
Reinforcement learning is gradient descent on a reward/uncertainty
signal --- the computational formalization of the same process Theorem 3
derives from survival pressure. RL's learning rate parameter η is the
computational analog of η\_ind in Theorem 3a.

\subsubsection{{[}Woolley2010{]}}\label{woolley2010}

Woolley, A. W., Chabris, C. F., Pentland, A., Hashmi, N., \& Malone, T.
W. (2010). Evidence for a collective intelligence factor in the
performance of human groups. \emph{Science}, 330(6004), 686--688.
https://doi.org/10.1126/science.1193147

\textbf{Supports:} Section 9 (Collective Knowledge), Theorem 6
(Collective Gradient Dominance), Corollary 15 (Self-Assembly).
\textbf{Stance:} Empirical support. Collective intelligence is
measurable and distinct from individual intelligence --- consistent with
Definition 15 (Superintelligence as a property of K\_collective, not any
individual E\_i).

\subsubsection{{[}Liu2024{]}}\label{liu2024}

Liu, N. F., Lin, K., Hewitt, J., Paranjape, A., Bevilacqua, M., Petroni,
F., \& Liang, P. (2024). Lost in the middle: How language models use
long contexts. \emph{Transactions of the Association for Computational
Linguistics}, 12, 157--173. https://doi.org/10.1162/tacl\_a\_00638

\textbf{Supports:} Corollary 4b (Scaling Argument Against Context Window
Growth), Discussion section (civilization-scale knowledge system).
\textbf{Stance:} Empirical support. Language models with longer contexts
perform worse at retrieving relevant information from non-salient
positions. This is the empirical signature of ρ degradation under
Theorem 2a: more context does not improve inference; it introduces noise
that degrades retrieval of the relevant subgraph K\^{}\{{[}f{]}\}.

\subsubsection{{[}vanderAalst2016{]}}\label{vanderaalst2016}

van der Aalst, W. M. P. (2016). \emph{Process Mining: Data Science in
Action} (2nd ed.). Springer. https://doi.org/10.1007/978-3-662-49851-4

\textbf{Supports:} Section 11 Related Work (Theorem 9, Observational
Bootstrapping). \textbf{Stance:} Adjacent method. Process mining
discovers workflow models from event logs --- the closest existing CS
method to equivalence class skill precipitation. Distinction: produces
flat workflow descriptions without provenance, confidence measures, or
connection to a transposability threshold.

\subsubsection{{[}AbbeelNg2004{]}}\label{abbeelng2004}

Abbeel, P., \& Ng, A. Y. (2004). Apprenticeship learning via inverse
reinforcement learning. \emph{Proceedings of the 21st International
Conference on Machine Learning (ICML)}, 1.
https://doi.org/10.1145/1015330.1015430

\textbf{Supports:} Section 11 Related Work (Theorem 9, Observational
Bootstrapping). \textbf{Stance:} Adjacent method. IRL infers reward
functions from demonstrated behavior; the output is a policy to imitate.
Distinction: infers what to optimize for, not an informed action with
grounded evidence, conditions, and provenance.

\subsubsection{{[}AamodtPlaza1994{]}}\label{aamodtplaza1994}

Aamodt, A., \& Plaza, E. (1994). Case-based reasoning: Foundational
issues, methodological variations, and system approaches. \emph{AI
Communications}, 7(1), 39--59.

\textbf{Supports:} Section 11 Related Work (Theorem 9, Observational
Bootstrapping). \textbf{Stance:} Adjacent method. CBR retrieves similar
past cases to inform new decisions. Distinction: retrieves the similar
case; Theorem 9 induces the general informed action from the equivalence
class, with formal confidence derived from \(k\) independent instances.

\subsubsection{{[}Muggleton1991{]}}\label{muggleton1991}

Muggleton, S. (1991). Inductive logic programming. \emph{New Generation
Computing}, 8(4), 295--318. https://doi.org/10.1007/BF03037089

\textbf{Supports:} Section 11 Related Work (Theorem 9, Observational
Bootstrapping). \textbf{Stance:} Adjacent method. ILP learns general
rules from positive and negative examples. Distinction: no provenance,
no uncertainty measure, no rigor threshold determining when the induced
rule is reliable enough to promote to lower encoding levels.

\subsubsection{{[}FinnEtAl2017{]}}\label{finnetal2017}

Finn, C., Abbeel, P., \& Levine, S. (2017). Model-agnostic meta-learning
for fast adaptation of deep networks. \emph{Proceedings of the 34th
International Conference on Machine Learning (ICML)}, 70, 1126--1135.

\textbf{Supports:} Section 11 Related Work (Theorem 9, Observational
Bootstrapping); Definition 11a (η\_M at collective scale).
\textbf{Stance:} Correspondence and extension. MAML learns an
initialization of M that adapts quickly across a task distribution ---
structurally η\_M at collective scale. Distinction: optimizes adaptation
speed; this framework additionally derives the rigor threshold for
promotion and the provenance requirement for transposability.

\begin{center}\rule{0.5\linewidth}{0.5pt}\end{center}

\subsection{Citation placement map}\label{citation-placement-map}

The following lists where each key needs to be inserted in THEORY.md.
Work through this list when adding inline markers.

{\def\LTcaptype{none} % do not increment counter
\begin{longtable}[]{@{}
  >{\raggedright\arraybackslash}p{(\linewidth - 4\tabcolsep) * \real{0.3333}}
  >{\raggedright\arraybackslash}p{(\linewidth - 4\tabcolsep) * \real{0.3333}}
  >{\raggedright\arraybackslash}p{(\linewidth - 4\tabcolsep) * \real{0.3333}}@{}}
\toprule\noalign{}
\begin{minipage}[b]{\linewidth}\raggedright
Location
\end{minipage} & \begin{minipage}[b]{\linewidth}\raggedright
Citation(s)
\end{minipage} & \begin{minipage}[b]{\linewidth}\raggedright
Note
\end{minipage} \\
\midrule\noalign{}
\endhead
\bottomrule\noalign{}
\endlastfoot
Abstract para 3 (FEP critique) & {[}Friston2010{]}, {[}Friston2017{]} &
Name FEP explicitly \\
Section 1, opening & {[}Friston2010{]} & ``departure from'' FEP \\
Definition 5 (H(w\textbar K)) & {[}Shannon1948{]}, {[}Cover2006{]},
{[}Pearl1988{]} & Grounding U formally; Bayesian network
interpretation \\
Definition 3b (I(f;w\textbar K)) & {[}Shannon1948{]} & Mutual
information \\
Lemma 1 proof (monotonicity) & {[}Cover2006{]} & ``conditional entropy
is monotonically non-increasing'' \\
Definition 3b G≥0 (data processing) & {[}Cover2006{]} & Data processing
inequality \\
Section 3 (K/F Inseparability) & {[}Gibson1979{]}, {[}Tishby2011{]} &
Affordances / bottleneck \\
Definition 6 (Agency A) & {[}Schmidhuber2010{]} & Formal curiosity
drive \\
Definition 7 (I(E) = A·η\_M) & {[}Schmidhuber2010{]}, {[}Thrun1998{]} &
Compression progress / meta-learning \\
Definition 8 / Theorem 2 (Escalation) & {[}Baars1988{]},
{[}Dehaene2011{]}, {[}Kahneman2011{]} & GWT / System 1-2 \\
Theorem 2a note (FEP self-undermining) & {[}Friston2010{]},
{[}Friston2017{]} & Precision allocation \\
Corollary 4b (context window scaling) & {[}Liu2024{]} & Empirical
support \\
Definition 11 / 11a (M, η\_M) & {[}Thrun1998{]} & Meta-learning \\
Definition 14 note (sentience disclaimer) & {[}Chalmers1995{]} & Hard
problem \\
Definition 16 (Truth as limit) & {[}Peirce1877{]}, {[}Peirce1878{]} &
Pragmatist truth \\
Theorem 6 (Collective Gradient) & {[}Woolley2010{]} & Empirical
collective intelligence \\
Section 5 thermodynamic commentary & {[}Schrodinger1944{]} & Negative
entropy \\
Open Question 1 & {[}Prigogine1984{]}, {[}England2013{]},
{[}Schrodinger1944{]} & Dissipative structures \\
Open Question 7 (boundary of M) & {[}Hinton1987{]} & Baldwin effect / M
acting on M \\
Open Question 9 (DNA as encoding of M) & {[}Schrodinger1944{]},
{[}MaynardSmith1995{]}, {[}Hinton1987{]}, {[}England2013{]} & Evolution
/ encoding \\
Discussion (civilization-scale) & {[}Liu2024{]}, {[}Woolley2010{]} &
Empirical grounding \\
Section 11 Related Work (FEP) & All Tier 1 and 2 & Full engagement \\
Theorem 9 / Section 11 (CS-side) & {[}vanderAalst2016{]},
{[}AbbeelNg2004{]}, {[}AamodtPlaza1994{]}, {[}Muggleton1991{]},
{[}FinnEtAl2017{]} & Observational bootstrapping \\
\end{longtable}
}

\end{document}
